% Generated draft: requires mathematical and source review.
\documentclass[11pt]{article}
\usepackage[margin=25mm]{geometry}
\usepackage{fontspec}
\setmainfont{Latin Modern Roman}
\usepackage{amsmath,amssymb,mathtools,unicode-math}
\setmathfont{Latin Modern Math}
\usepackage{enumitem,needspace,xurl,hyperref}
\setlength{\parindent}{0pt}
\setlength{\parskip}{5pt}
\setlength{\emergencystretch}{3em}
\newcommand{\N}{\mathbb{N}}
\newcommand{\Z}{\mathbb{Z}}
\newcommand{\Q}{\mathbb{Q}}
\newcommand{\R}{\mathbb{R}}
\newcommand{\C}{\mathbb{C}}
\newcommand{\F}{\mathbb{F}}
\newcommand{\A}{\mathbb{A}}
\newcommand{\PP}{\mathbb{P}}
\newcommand{\E}{\mathbb{E}}
\newcommand{\eps}{\varepsilon}
\newcommand{\dd}{\,\mathrm{d}}
\begin{document}
\section*{UnsolvedMath ID: 30005856}
\textbf{Source problem code:} catalogue entry
\subsection{Definitions and mathematical statement}
Let \((M^n,g)\) be a smooth Riemannian manifold. Write \(R_g\) for its scalar curvature, \(\operatorname{Ric}_g\) for its Ricci tensor, and \(\operatorname{Rm}_g\) for its Riemann curvature tensor. Use the curvature-operator convention in which the round sphere has positive curvature operator. At each \(x\in M\), let
\[
\lambda_1(x)\leq\lambda_2(x)\leq\cdots\leq\lambda_{N}(x),\qquad N=\binom n2,
\]
be the eigenvalues, counted with multiplicity, of the self-adjoint curvature operator on \(\Lambda^2T_xM\). In particular, the expression involving two eigenvalues requires \(n\geq3\). The supplied record calls the metric Ricci-pinched and \(2\)-pinched when there is a constant \(c>0\), independent of \(x\), such that, everywhere on \(M\),
\[
R_g\geq0,\qquad \operatorname{Ric}_g\geq cR_g g,\qquad \lambda_1+\lambda_2\geq cR_g.
\]
The tensor inequality means that \(\operatorname{Ric}_g(v,v)\geq cR_g(x)g(v,v)\) for every \(x\in M\) and every \(v\in T_xM\). The constant may depend on the metric. Using one constant for both pinching inequalities is equivalent to allowing two positive constants, because \(R_g\geq0\) and their minimum can be used.

The intended question is whether these uniform pinching conditions and completeness force the dichotomy
\[
\operatorname{Rm}_g\equiv0\qquad\text{or}\qquad M\text{ is compact},
\]
without an additional noncollapsing hypothesis. Here flatness means \(\operatorname{Rm}_g\equiv0\), and completeness is understood as geodesic completeness under the usual convention for manifolds without boundary.

There are unresolved domain qualifications. The supplied reconstructed statement does not specify connectedness or whether manifolds have boundary. A proposed precise completion is to ask the displayed question for every complete, connected Riemannian manifold without boundary in each integer dimension \(n\geq4\), with \(c>0\) quantified separately for each metric. Connectedness and the boundary convention in this completion require confirmation from the source; they are not explicit hypotheses of the supplied reconstructed statement. Without connectedness, the global flat-or-compact dichotomy is not valid in general. The phrase higher dimensions suggests \(n\geq4\), but the exact intended dimension range is not explicitly supplied.

No bounded-curvature hypothesis is included in the reconstructed question. The quoted discussion mentions bounded curvature, connectedness, nonnegative Ricci curvature, and nonflatness for starting metrics in a dimension-three proof discussion; that passage does not establish these as additional hypotheses of the higher-dimensional question. Its references to existence and noncollapsing results likewise describe proof background rather than an assumed noncollapsing condition in the question.
\subsection{Short English statement}
The intended question asks whether a complete manifold with uniformly positive Ricci and two-eigenvalue pinching relative to its nonnegative scalar curvature must be flat or compact, without assuming noncollapse. The higher-dimensional range and the necessary connectedness and boundary conventions require source confirmation. The supplied status is partially solved; its dated literature assessment is not independently verified here.
\subsection{Sources}
\begin{itemize}
\item Supplied local UnsolvedMath registry record, permanent ID \(30005856\), title: Higher-Dimensional Hamilton-Lott Conjecture. Collection: Oberwolfach Reports Open Problems. Exact source problem code: \texttt{catalogue entry}; proposed year: 2024. No individual proposer is specified.
\item Supplied report citation: Analysis, Geometry and Topology of Positive Scalar Curvature Metrics (2024), Oberwolfach Reports, pp.\ 483--566, DOI \texttt{10.4171/owr/2024/9}. \url{https://doi.org/10.4171/owr/2024/9}. The background describes Oberwolfach Reports as documenting mathematical workshops and their open-problem discussions.
\item The supplied excerpt attributes the presentation On the Hamilton-Lott conjecture in higher dimensions to Alix Deruelle, joint work with Felix Schulze and Miles Simon. It locates the higher-dimensional question at reference [3, Question 1.5]. Complete bibliographic identifications of the excerpt's references [3] and [8] are not supplied.
\item Supplied literature source: On the Hamilton-Lott conjecture in higher dimensions. \url{https://arxiv.org/abs/2403.00708}.
\item Supplied literature source: Oberwolfach Report on the Hamilton-Lott conjecture. \url{https://ems.press/content/serial-article-files/48658?nt=1}.
\item Literature assessment dated exactly \texttt{2026-08-22}: the supplied curation labels the problem partially solved and reports that the Hamilton-Lott conjecture has been proved in dimension three and in higher-dimensional uniformly PIC1-pinched non-collapsed settings, while the broader higher-dimensional \(2\)-pinching question is not fully settled. These are supplied dated curation claims, not an independently verified literature assessment. The supplied background expressly requests independent verification of time-sensitive status and cited claims.
\item The registry labels the reconstructed statement \texttt{corrected\_verified} and states that it was reconstructed and checked against the official Oberwolfach report. This is a supplied provenance claim; the source has not been independently visited or checked for this draft.
\end{itemize}
\end{document}
